\par\noindent\textbf{(1)} 相邻两槽的入射光程差为 \(d\sin\alpha\)，出射方向光程差为 \(d\sin\beta\)。相长干涉需 \(k d(\sin\alpha+\sin\beta)=2\pi m\)，所以 \(d(\sin\alpha+\sin\beta)=m\lambda\)，只有满足 \(|m\lambda/d-\sin\alpha|\le1\) 的级次能传播。
\par\noindent\textbf{(2)} 对光栅方程固定 \(\alpha\) 求波长微分，得 \(d\cos\beta\,d\beta=m,d\lambda\)，因此 \(d\beta/d\lambda=m/(d\cos\beta)\)。当 \(\cos\beta\) 很小时形式角度色散很大，但掠射角光学效率及实际孔径变得关键。
\par\noindent\textbf{(3)} 对 \(N\) 条等间距槽，相位因子之和形成 \(\sin(N\Delta/2)/\sin(\Delta/2)\) 的主峰；第一零点相位偏移约 \(2\pi/N\)。将两个邻近谱线分离到Rayleigh极限，得 \(\lambda/\Delta\lambda\approx mN\)，不是单独由槽宽决定。
\par\noindent\textbf{(4)} 若锯齿槽面法线相对总法线倾角为 \(\theta_B\)，单个斜面镜面反射最强方向近似满足 \(\alpha+\beta=2\theta_B\)（按约定同侧角）。同时需满足光栅衍射方程，才能使所选级次落在闪耀包络峰；真实效率还需考虑偏振和槽形矢量衍射。
